% drafts/06-exactness.tex — Algebraic Exactness of the Pauli Pipeline
%
% Target length: ~1.5 pages (two-column PRA)
% Subsections: 6.1 Phase representation, 6.2 Pauli multiplication,
%              6.3 Like-term combination, 6.4 Exactness theorem

Every fermion-to-qubit encoding ultimately reduces to multiplying and
summing Pauli strings with complex coefficients.  In standard
implementations (e.g., OpenFermion~\cite{mcclean2020},
PennyLane~\cite{bergholm2022}), phases such as $\pm 1$ and $\pm i$
are represented as floating-point complex numbers.  For single products
this introduces no visible error, but in deep multiplication chains
and Hamiltonian assembly with thousands of terms, floating-point
cancellation can accumulate silently.  We show that the Pauli pipeline
can be made \emph{exactly} symbolic by tracking phases in the cyclic
group $\Zi$.

%%====================================================================
\subsection{The phase representation}
\label{sec:ex-phase}

The four values $\{+1, -1, +i, -i\}$ arising from Pauli products form
the cyclic group $\Zi$ under multiplication (\cref{def:phase-group}).  We
represent this group as a four-element discriminated union:
\begin{align}
  \texttt{P1} &\leftrightarrow +1, &
  \texttt{M1} &\leftrightarrow -1, \notag \\
  \texttt{Pi} &\leftrightarrow +i, &
  \texttt{Mi} &\leftrightarrow -i.
  \label{eq:phase-enum}
\end{align}
Multiplication is implemented by exhaustive pattern matching on the
$4 \times 4$ Cayley table.  No floating-point operation is involved.

The key operation \textsc{FoldIntoGlobalPhase} converts a symbolic
phase into an effect on a \texttt{Complex} coefficient using only
negation and the swap $(a + bi) \mapsto (-b + ai)$ for multiplication
by~$i$.  This avoids the standard \texttt{Complex.Multiply} call and
its associated rounding.

%%====================================================================
\subsection{Pauli multiplication}
\label{sec:ex-pauli}

The product of two single-qubit Pauli operators returns a pair
$(\sigma_k, \phi)$ where $\sigma_k \in \{I, X, Y, Z\}$ and
$\phi \in \Zi$:
\begin{equation}
  \sigma_a \cdot \sigma_b = \phi_{ab}\; \sigma_{a \oplus b},
  \label{eq:pauli-product}
\end{equation}
where $\oplus$ is the mod-2 group operation on the Pauli index
labels.  The full table is:
\begin{center}
\begin{tabular}{@{}c|cccc@{}}
$\cdot$ & $I$ & $X$ & $Y$ & $Z$ \\
\midrule
$I$ & $(I,+1)$ & $(X,+1)$ & $(Y,+1)$ & $(Z,+1)$ \\
$X$ & $(X,+1)$ & $(I,+1)$ & $(Z,+i)$ & $(Y,-i)$ \\
$Y$ & $(Y,+1)$ & $(Z,-i)$ & $(I,+1)$ & $(X,+i)$ \\
$Z$ & $(Z,+1)$ & $(Y,+i)$ & $(X,-i)$ & $(I,+1)$ \\
\end{tabular}
\end{center}
This is computed by pattern matching with no arithmetic.

For an $n$-qubit Pauli string, multiplication proceeds qubit-by-qubit
in $O(n)$ time.  At each position, the single-qubit product returns
$(\sigma_k, \phi_k)$.  The phase factors are accumulated via the
$\Zi$ group law, then folded into the global coefficient exactly once
at the end.  The coefficient $c_\text{global}$ is modified only by
the \textsc{FoldIntoGlobalPhase} operation — never by floating-point
complex multiplication of phase factors.

%%====================================================================
\subsection{Like-term combination}
\label{sec:ex-combination}

After encoding a Hamiltonian, distinct products may map to the same
Pauli string.  Combining like terms is done by a dictionary keyed on
the string \emph{signature} — the operator labels without coefficients
(e.g., ``$XYZII$'').

\begin{algorithm}[t]
\caption{Like-term combination}
\label{alg:like-term}
\begin{algorithmic}[1]
\State $D \gets$ empty dictionary
\For{each Pauli register $r$ with signature $s$ and coefficient $c$}
  \If{$s \in D$}
    \State $c' \gets D[s].\text{coefficient} + c$
    \If{$c' = 0$}
      \State remove $s$ from $D$ \Comment{exact cancellation}
    \Else
      \State $D[s] \gets r$ with coefficient $c'$
    \EndIf
  \Else
    \State $D[s] \gets r$
  \EndIf
\EndFor
\end{algorithmic}
\end{algorithm}

Because the dictionary key is a \emph{string} (not a floating-point
comparison), matching is exact.  Two terms contribute to the same
entry if and only if they act on precisely the same qubits with
precisely the same Pauli operators.  Coefficient addition is ordinary
complex addition — the only point where floating-point arithmetic
enters — and cancellation (removal of zero-coefficient terms) is
checked against exact zero.

%%====================================================================
\subsection{The exactness theorem}
\label{sec:ex-theorem}

\begin{theorem}[Algebraic exactness]
\label{thm:exactness}
Let $H = \sum_k c_k \prod_j O_j^{(k)}$ be a fermionic Hamiltonian
with coefficients $c_k \in \Qi$.  The Pauli-string representation
$\encode(H) = \sum_m d_m P_m$ produced by the encoding pipeline
satisfies $d_m \in \Qi$ for all~$m$, and the computation introduces
no floating-point error in the phase tracking.
\end{theorem}

\begin{proof}[Proof sketch]
Each encoding step produces Majorana operators whose coefficients are
$\pm 1$ or $\pm i$ (elements of $\Zi$), combined with the
input coefficients $c_k$.  The encoding pipeline consists of:
\begin{enumerate}
  \item \emph{Ladder $\to$ Majorana}: coefficients are
    $\pm\frac{1}{2}$ and $\pm\frac{i}{2}$ (elements of
    $\frac{1}{2}\Qi$).
  \item \emph{Majorana $\to$ Pauli string}: each Majorana operator is
    a single Pauli string with coefficient $\pm 1$ or $\pm i$,
    tracked symbolically in $\Zi$.
  \item \emph{Pauli $\times$ Pauli}: qubit-by-qubit multiplication
    via the finite table of \cref{eq:pauli-product}, phase
    accumulated in $\Zi$.
  \item \emph{Phase fold}: a single application of
    \textsc{FoldIntoGlobalPhase}, using only negation and
    real/imaginary swap.
  \item \emph{Like-term combination}: coefficient addition in
    $\mathbb{C}$, with exact string matching.
\end{enumerate}
Steps 1--4 introduce no rounding.  Step 5 performs
$\Qi$-coefficient additions that are exact in IEEE~754 when
the operands are half-integers (representable exactly as
\texttt{double}).  The full proof is in \cref{app:exactness-proof}.
\end{proof}

\begin{remark}
When input coefficients are not in $\Qi$ (e.g., from quantum
chemistry integrals), the pipeline still benefits from exact phase
tracking: the only source of rounding is the input data itself,
not the encoding algebra.  In contrast, libraries that multiply
phases as \texttt{Complex(0, 1)} accumulate $O(\epsilon)$ errors
per Pauli product, where $\epsilon \approx 10^{-16}$ is machine
epsilon.  Over $K$ products, the worst-case phase error grows as
$K\epsilon$, which can become significant for large Hamiltonians.
\end{remark}
